% Краткое сообщение для УМН. Оформление --- официальный шаблон
% (mathnet.ru/rus/rm/authornotes, umn_shap.zip: 000.tex + umnbib.sty),
% перекодированный в UTF-8. Формальный предел жанра --- две журнальные
% страницы; эта редакция объявляет все пять оценок и потому длиннее.
% Двухстраничный вариант без результата в R^10 --- note-2p.tex (см. README.md).
\documentclass{article}
\usepackage[T2A]{fontenc}
\usepackage[utf8]{inputenc}
\usepackage[russian]{babel}
\usepackage[tbtags]{amsmath}
\usepackage{amsfonts,amssymb,mathrsfs,amscd,comment}
\usepackage{umnbib} % оформление списка литературы

\overfullrule10pt

\voffset-30mm\hoffset-15mm\mag1200
\textheight 230mm\textwidth 140mm\normalbaselineskip=12.5pt %4.39mm
\setlength{\emergencystretch}{1.5em}

\newcommand{\R}{\mathbb{R}}
\newcommand{\Z}{\mathbb{Z}}
\newcommand{\CC}{\mathbb{C}}
\newcommand{\diam}{\operatorname{diam}}
\newcommand{\dist}{\operatorname{dist}}
\newtheorem{theorem}{Теорема}
\newtheorem{lemma}{Лемма}
\newtheorem{proposition}{Предложение}

\begin{document}

% TODO: сверить УДК с редакцией
\title{\makebox[\textwidth][l]{\normalfont\normalsize УДК 514.174.5+519.174}\\[8pt]
\bf\large\MakeUppercase{%
Новые верхние оценки хроматических чисел евклидовых пространств
}}

% Порядок авторов --- по вкладу (решение от 22.08.2026): в заметку входят
% решётка индекса 43 и точные верификаторы 1029 и 7203 (Н.\,В.~Глушкова).
\author{Л.\,Л.~Иванов, Н.\,В.~Глушкова}
\date{}

\maketitle

\makeatletter
\renewcommand{\@makefnmark}{}
\makeatother

\textbf{1. Постановка и результат.}
Через $\chi(\R^n,[1,\ell])$, $\ell\ge1$, обозначим наименьшее число цветов
раскраски пространства $\R^n$, при которой никакие две одноцветные точки не
находятся на расстоянии из отрезка $[1,\ell]$; при $\ell=1$ это классическое
хроматическое число $\chi(\R^n)$ (проблема Нельсона--Хадвигера, см.~\cite{Soifer});
интервальная постановка изучалась в~\cite{Ivanov2011}. Очевидно,
$\chi(\R^n)\le\chi(\R^n,[1,\ell])$. Лучшие известные верхние оценки в
размерностях $4$, $5$, $7$, $9$, $10$ --- соответственно $49$, $140$, $1372$,
$17253$ и $3^{10}=59049$ \cite{ABPR}; в \cite{ABPR} высказано ожидание, что
$49$ и $140$ неулучшаемы решёточными раскрасками.

\begin{theorem}
$\chi(\R^4)\le43$, $\chi(\R^5)\le132$, $\chi(\R^7)\le1029$,
$\chi(\R^9)\le7203$, $\chi(\R^{10})\le45619$. Точнее,
$\chi(\R^n,[1,\ell])\le k$ при $(n,k)=(4,43)$ для $\ell\le1{,}00411$,
при $(5,132)$ для $\ell\le101/100$, при $(7,1029)$ для $\ell\le103/100$,
при $(9,7203)$ для $\ell<1{,}0166127$ и при $(10,45619)$ для
$\ell<\sqrt{217/214}$.
\end{theorem}

Оценка $45619$ --- первая в $\R^{10}$ ниже классической величины $3^n$.
Первые четыре оценки даны явными рациональными решётками с проверкой в
точной арифметике (п.~3), последняя выведена аналитически (п.~4); подробности
--- \cite{Full}.

\textbf{2. Решёточные раскраски.}
Пусть $\Lambda\subset\R^n$ --- решётка, $\Gamma\subset\Lambda$ --- подрешётка
индекса $k$, $V_0$ --- ячейка Вороного нуля (выпуклый центрально-симметричный
многогранник); ячейку $a+V_0$, $a\in\Lambda$, окрасим цветом класса $a+\Gamma$.
Положим
\[
D(v)=2\dist\bigl(v/2,\,V_0\bigr),\qquad
D(\Gamma)=\min_{v\in\Gamma\setminus\{0\}}D(v),\qquad
d(\Lambda,\Gamma)=\frac{D(\Gamma)}{\diam V_0}.
\]
Если $x\in V_0$, $y\in v+V_0$, то $x'=v-y\in V_0$ и
$x-y=2\bigl(\tfrac{x+x'}2-\tfrac v2\bigr)$, откуда по выпуклости
$|x-y|\ge D(v)$ (на самом деле равенство, \cite[лемма~2]{Ivanov2011}). Значит,
две одноцветные точки либо лежат в одной ячейке (расстояние $\le\diam V_0$),
либо удалены не менее чем на $D(\Gamma)$. Если $d=d(\Lambda,\Gamma)>1$, то при
масштабе $s$ с $s\diam V_0<1$, $sD(\Gamma)>\ell$ (такой $s$ есть в точности при
$\ell<d$) отрезок $[1,\ell]$ свободен от одноцветных расстояний, то есть
\begin{equation}\label{crit}
\chi(\R^n,[1,\ell])\le[\Lambda:\Gamma]\quad\text{при всех }\ell<d(\Lambda,\Gamma).
\end{equation}
Величину $d$ будем называть \emph{шириной} раскраски.

\textbf{3. Четыре явные конструкции.}
Пусть $Q$ --- рациональная положительно определённая матрица, $\Lambda=\Z^n$
со скалярным произведением $\langle x,y\rangle=x^{\mathsf T}Qy$ и
$\Gamma\subset\Lambda$ --- целочисленная подрешётка. Тогда неравенство
$d(\Lambda,\Gamma)>\ell_0$ сводится к конечной проверке в рациональных числах:
фасеты $V_0$ задаются релевантными векторами (по теореме Вороного это
единственные пары $\pm v$ кратчайших представителей ненулевых классов
$\Lambda/2\Lambda$), вершины $V_0$ --- решения рациональных линейных систем,
$\diam V_0=2R$, где $R^2$ --- максимум $\langle x,x\rangle$ по вершинам;
нарушить $D(v)\ge\ell_0\diam V_0$ может лишь вектор с $|v|<2(1+\ell_0)R$
(ибо $D(v)\ge|v|-\diam V_0$), таких конечное число, и для каждого
$\dist(v/2,V_0)$ подтверждается условиями Каруша--Куна--Таккера в дробях.
Для четырёх решёток ниже проверка выполнена целиком в точной арифметике;
полные данные --- в \cite{Full}.

\emph{$n=4$, $k=43$.} Решётка эйзенштейнова: $\Lambda=\Z[\omega]^2$,
$\omega=e^{2\pi i/3}$, с эрмитовой формой
$\bigl(\begin{smallmatrix}a&c\\\bar c&b\end{smallmatrix}\bigr)$, вещественная
часть которой в $\Z$-базисе $(e_1,\omega e_1,e_2,\omega e_2)$ при
$u=\operatorname{Re}c$, $s=\tfrac{\sqrt3}2\operatorname{Im}c$ имеет матрицу
Грама
\[
Q=\begin{pmatrix}
a & -\tfrac a2 & u & -\tfrac u2-s\\
-\tfrac a2 & a & -\tfrac u2+s & u\\
u & -\tfrac u2+s & b & -\tfrac b2\\
-\tfrac u2-s & u & -\tfrac b2 & b
\end{pmatrix},\quad
a=1,\ b=\tfrac{39261}{25000},\ u=\tfrac{37381}{100000},\ s=\tfrac{3141}{12500};
\]
подрешётка $\Gamma$ порождена строками матрицы
$T=\Bigl(\begin{smallmatrix}1&0&0&41\\0&1&0&14\\0&0&1&37\\0&0&0&43\end{smallmatrix}\Bigr)$,
$\det T=43$ ($43=N(7+\omega)$, $\Gamma$ --- $\Z[\omega]$-подмодуль). Ячейка
имеет $30$ фасет и $120$ вершин,
\[
\diam^2V_0=\tfrac{5374343884337}{2019775849050},\quad
D(\Gamma)^2=\tfrac{7396075258066147}{2756867811550000},\quad
d=1{,}0041105\ldots>1{,}00411 .
\]

\emph{$n=5$, $k=132$.} Решётка общего положения:
\[
Q=\tfrac1{100000}\begin{pmatrix}
79327&86284&69597&36140&26501\\
86284&254538&216808&101257&102135\\
69597&216808&296568&163667&140527\\
36140&101257&163667&181733&71653\\
26501&102135&140527&71653&148076
\end{pmatrix},
\quad
\Gamma=\ker\varphi,
\]
$\varphi(x)=88x_1+89x_2+127x_3+57x_4+3x_5 \bmod 132$ (гомоморфизм сюръективен,
$[\Lambda:\Gamma]=132$). Ячейка имеет $62$ фасеты и $720$ вершин, в окне
$|v|<2(1+\tfrac{101}{100})R$ лежат $38$ ненулевых векторов $\Gamma$,
{\small
\[
R^2=\tfrac{3093570095915736710467707313641}{3999974238630281363046165200000},\qquad
D(\Gamma)^2=\tfrac{1634114770527727}{516898614620000},\qquad
d=1{,}0108977\ldots>\tfrac{101}{100}.
\]}

\emph{$n=7$, $k=1029$.} Ламинирование эйзенштейновой раскраски
$(3+\omega)E_6^*$ индекса $7^3=343$ тремя слоями: в подходящем $\Z$-базисе
\[
Q=\tfrac1{10^4}\begin{pmatrix}
30000&15000&0&-22500&15000&7500&-8632\\
15000&30000&22500&0&7500&15000&-4197\\
0&22500&45000&22500&0&22500&1601\\
-22500&0&22500&45000&-22500&0&6617\\
15000&7500&0&-22500&30000&15000&-5275\\
7500&15000&22500&0&15000&30000&4434\\
-8632&-4197&1601&6617&-5275&4434&17897
\end{pmatrix},
\quad
\Gamma=C\Z^7,
\]
где $C$ --- верхнетреугольная матрица с диагональю $(7,1,7,1,7,1,3)$ и
единственными ненулевыми внедиагональными элементами
$C_{12}=C_{34}=C_{56}=-5$; $\det C=1029$. Ячейка имеет $254$ фасеты и
$30\,368$ вершин,
\[
\diam^2V_0=\tfrac{96858928789031597}{14761819462500000},\qquad
D(\Gamma)^2=7,\qquad
d=1{,}0328782\ldots>\tfrac{103}{100};
\]
минимум $D(\Gamma)^2=7$ --- планарная граница предложения~1 ниже.

\emph{$n=9$, $k=7203$.} То же ламинирование над $(3+\omega)E_8$ индекса
$7^4=2401$ с тремя слоями высоты $t^2=\tfrac{8264707}{6250000}$: $Q=G/10^4$,
$\Gamma=C\Z^9$, $|\det C|=7203=3\cdot7^4$ (явные $G$, $C$ --- в \cite{Full}).
Ячейка имеет $752$ фасеты и $1\,654\,230$ вершин, из них $4590$ непростых,
\[
\diam^2V_0=\tfrac{1119552292542693}{165294140000000},\qquad
D(\Gamma)^2=7,\qquad d=1{,}0166127\ldots>1 .
\]
Априорной границы $\diam^2V_0\le6+t^2$ не хватает: нужен полный список вершин,
и его полнота доказана замыканием по $1$-скелету без предположения простоты.

В каждом из четырёх случаев \eqref{crit} даёт утверждение теоремы; для первых
трёх проверка независимо повторена второй программой без общих геометрических
процедур.

\textbf{4. Продуктовое правило и размерность $10$.}

\begin{lemma}
Пусть $\Lambda=\bigoplus_{i=1}^{s}\Lambda_i$ --- ортогональная прямая сумма
решёток $\Lambda_i\subset\R^{n_i}$, $\Gamma=\bigoplus_i\Gamma_i$,
$\Gamma_i\subseteq\Lambda_i$, $k_i=[\Lambda_i:\Gamma_i]$, и пусть
$d(\Lambda_i,\Gamma_i)\ge d_i>0$. Тогда при подходящем выборе масштабов
блоков $d(\Lambda,\Gamma)\ge\bigl(\sum_{i}d_i^{-2}\bigr)^{-1/2}$, и потому
$\chi(\R^{n},[1,\ell])\le\prod_i k_i$ при
$\ell<\bigl(\sum_{i}d_i^{-2}\bigr)^{-1/2}$, где $n=\sum_i n_i$.
\end{lemma}

\textbf{Доказательство.}
Масштабирование блока $(\Lambda_i,\Gamma_i)$ не меняет $d(\Lambda_i,\Gamma_i)$, так что
величины $s_i=\diam^2V_0(\Lambda_i)$ можно задавать произвольно. Ячейка
Вороного прямой суммы есть произведение ячеек, поэтому
$\diam^2V_0(\Lambda)=\sum_i s_i$, а для $v=(v_1,\dots,v_s)$
$\dist^2(v/2,V_0)=\sum_i\dist^2\bigl(v_i/2,V_0(\Lambda_i)\bigr)$, то есть
$D(v)^2=\sum_iD_i(v_i)^2$. Так как $D_i(0)=0$, минимум по
$\Gamma\setminus\{0\}$ достигается на векторе с единственной ненулевой
компонентой: $D(\Gamma)^2=\min_iD(\Gamma_i)^2\ge\min_i d_i^2s_i$. Положив
$s_i=S/d_i^2$, получаем $D(\Gamma)^2\ge S$ и
$\diam^2V_0(\Lambda)=S\sum_id_i^{-2}$, откуда
$d(\Lambda,\Gamma)^2\ge\bigl(\sum_id_i^{-2}\bigr)^{-1}$. Индекс $\Gamma$ в
$\Lambda$ равен $\prod_ik_i$; остаётся применить \eqref{crit}. $\square$

Правая часть монотонна по $d_i$, поэтому достаточно оценок ширин блоков
\emph{снизу}; величину $d_i^{-2}$ назовём \emph{ценой} блока.

Решётку $\Lambda$ назовём \emph{эйзенштейновой}, если она снабжена
структурой модуля над $\Z[\omega]$, причём умножение на $\omega$ ---
изометрия. Пусть $H_1\subset\CC$ --- ячейка Вороного решётки
$\Z[\omega]\subset\CC=\R^2$: правильный шестиугольник с фасетами на расстоянии
$1/2$ от нуля и вершинами на расстоянии $1/\sqrt3$; $\lambda_1(\Lambda)$ ---
длина кратчайшего ненулевого вектора.

\begin{proposition}
Пусть $\Lambda$ эйзенштейнова и $\alpha\in\Z[\omega]\setminus\{0\}$. Тогда
$D(\alpha\Lambda)\ge2\,\lambda_1(\Lambda)\cdot\dist(\alpha/2,\,H_1)$.
\end{proposition}

\textbf{Доказательство.}
Всякий ненулевой вектор $\alpha\Lambda$ имеет вид $\alpha w$, $w\in\Lambda
\setminus\{0\}$. Векторы $\pm w$, $\pm\omega w$, $\pm(1+\omega)w=\mp\omega^2w$
лежат в $\Lambda$ и имеют длину $|w|$. Шесть полупространств
$\{x:\langle x,u\rangle\le|u|^2/2\}$ по этим $u$ содержат $V_0$, а их нормали
лежат в плоскости $P=\operatorname{span}(w,\omega w)$; поэтому ортогональная
проекция $V_0$ на $P$ содержится в шестиугольнике $H$, высекаемом в $P$ теми
же неравенствами, --- ячейке Вороного плоской решётки $\Z[\omega]w\subset P$,
изометричной $|w|\,H_1$. Точка $\alpha w/2$ лежит в $P$, и проекция не
увеличивает расстояний, значит,
$\dist(\alpha w/2,V_0)\ge\dist_P(\alpha w/2,H)=|w|\dist(\alpha/2,H_1)$. Так как
$|w|\ge\lambda_1(\Lambda)$, получаем требуемое. $\square$

Нужны два значения: ближайшая к точке $(3+\omega)/2=(5/4,\sqrt3/4)$ точка
$H_1$ --- вершина $(1/2,\,1/(2\sqrt3))$, квадрат расстояния $9/16+1/48=7/12$;
для $(5+2\omega)/2=(2,\sqrt3/2)$ ближайшая та же вершина, квадрат расстояния
$9/4+1/3=31/12$.

\emph{Блок $E_8$.} Решётка $E_8$ эйзенштейнова (см.~\cite{ABPR},
\cite{ConwaySloane}), $\lambda_1(E_8)^2=2$, её радиус покрытия равен~$1$
\cite[гл.~4]{ConwaySloane}, так что $\diam V_0(E_8)=2$. Подрешётка
$\Gamma_1=(3+\omega)E_8$ имеет индекс $N(3+\omega)^4=7^4=2401$
(конструкция из~\cite{ABPR}), и по предложению~1
$D(\Gamma_1)\ge2\sqrt2\cdot\sqrt{7/12}=\sqrt{14/3}$. Отсюда
$d(E_8,\Gamma_1)\ge\sqrt{14/3}\big/2=\sqrt{7/6}$, то есть цена блока не
превосходит $6/7$.

\emph{Плоский блок.} $\Lambda_2=\Z[\omega]\subset\CC$ (шестиугольная решётка,
$\lambda_1=1$, $\diam V_0=2/\sqrt3$), $\Gamma_2=(5+2\omega)\Z[\omega]$, индекс
$N(5+2\omega)=19$. По предложению~1 $D(\Gamma_2)\ge2\sqrt{31/12}$, поэтому
$d_2\ge\sqrt{31/12}\cdot\sqrt3=\sqrt{31/4}$ и цена не превосходит $4/31$.

Применим лемму к $E_8\oplus a\Z[\omega]$ с
$\Gamma_1\oplus(5+2\omega)a\Z[\omega]$ при подходящем $a>0$: так как
$\frac67+\frac4{31}=\frac{214}{217}<1$, получаем
$\chi(\R^{10},[1,\ell])\le2401\cdot19=45619$ при $\ell<\sqrt{217/214}$.
$\square$

\begin{thebibliography}{9}
\bibitem{Soifer} A.~Soifer, \emph{The Mathematical Coloring Book}, Springer,
New York, 2009.
\bibitem{Ivanov2011} Л.~Л.~Иванов, \emph{О хроматических числах $\R^2$ и $\R^3$
с интервалами запрещённых расстояний}, Вестник РУДН, сер. Матем. Информ. Физ.,
2011, №~1, 14--29.
\bibitem{ABPR} A.~Arman, A.~Bondarenko, A.~Prymak, D.~Radchenko, \emph{Upper
bounds on chromatic number of $\mathbb E^n$ in low dimensions}, Electron. J.
Combin., \textbf{31}:2 (2024), P2.35.
\bibitem{ConwaySloane} J.~H.~Conway, N.~J.~A.~Sloane, \emph{Sphere Packings,
Lattices and Groups}, 3rd ed., Springer, 1999.
\bibitem{Full} Л.~Л.~Иванов, Н.~В.~Глушкова, \emph{Новые верхние оценки
хроматических чисел евклидовых пространств: полная версия} (электронное
дополнение), 2026, \texttt{https://github.com/ivaleo/Chromatic/tree/main/paper}.
\end{thebibliography}

% TODO: место работы обоих авторов (для второго автора --- вместо <<---->>)
\noindent
\parbox[t]{0.5\textwidth}{{\bf Л.\,Л.~Иванов (L.\,L.~Ivanov)}\\
Москва\\
{\it E-mail}: leo.ivanov@gmail.com}%
\parbox[t]{0.5\textwidth}{{\bf Н.\,В.~Глушкова (N.\,V.~Glushkova)}\\
---\\
{\it E-mail}: glushkova.nv@phystech.edu}

\end{document}
